\section{Introduction}

The web has become the default training corpus for identity-hungry generative models, and faces are its most sensitive commodity. A handful of public portraits suffices to synthesize a person in any context: DreamBooth-style fine-tuning binds a rare token to an identity with four to eight images~\cite{ruiz2022dreambooth}, while zero-shot adapters such as InstantID and IP-Adapter-FaceID inject a face embedding from a single reference image directly into the denoising process~\cite{wang2024instantid,ye2023ipadapter}. The consent asymmetry is severe---the subject never agreed---and it has triggered an arms race in \emph{protective perturbation}: cloaks added to a photo before upload that break the customization pipeline. Anti-DreamBooth established the category~\cite{le2023antidreambooth}, and a fast-growing successor literature hardens it along every laboratory axis~\cite{liu2023metacloak,zheng2023targeted,onikubo2024highfrequency,li2024promptindependent,song2024idprotector}. The narrative is reassuring: with the right perturbation, your face is yours again.

\textbf{Relevance to the Web.} Every component of this study is Web infrastructure: the protection gap is created and erased by platform image pipelines (CDNs, avatar services, thumbnailers), browser rendering, and crawling agents; the adversary is any scraper harvesting faces the way the Web serves them; and the proposed remedy is a machine-readable consent manifest served over the Web's ingestion layer. The paper answers a Responsible Web research challenge---measuring, then repairing, the distance between the privacy users believe they control and the acquisition pipeline the Web actually executes.

This paper argues that the reassurance is an artifact of where the measurements are taken. Every published protection is evaluated on \emph{T0-images}: clean, full-resolution, uncompressed files that the adversary kindly provides to the defense. That world does not exist on the web. An identity that is ``public'' has already been transformed before a scraper ever sees it: platforms re-encode uploads under storage and bandwidth budgets, browsers rasterize images under device pixel ratios and thumbnail pipelines, crawling agents screenshot pages and crop detected faces rather than download files, and users regenerate their photos through diffusion filters before posting. Each transformation is individually harmless to recognition; jointly, they describe the actual data distribution any real adversary can harvest. We call this pipeline the \emph{web acquisition funnel} and the resulting blind spot the \emph{Web Acquisition Gap}: the difference between the protection a defense delivers on the files it was handed and the protection that survives on the files an adversary can actually obtain---an attack surface that requires no adversarial optimization, costs milliseconds to seconds per image, and operates \emph{before} the defense is ever consulted.

Two observations sharpen the claim. First, the funnel is not hypothetical: a single JPEG re-encode at quality 85--95 neutralizes a broad class of protection perturbations~\cite{sandovalsegura2023jpeg}, common image transformations systematically suppress unlearnable-example attacks~\cite{qin2023learning,sandovalsegura2023what}, and robustness under diffusion-aware adversaries remains unfavorable~\cite{peng2025contrastive}. Nor is the funnel's parameter range speculative: our probes of production image pipelines (Section~\ref{sec:reality}) land squarely inside the transcode tiers we evaluate. What is missing is a systematic, tiered, budget-aware account of the entire funnel under both tunable and zero-shot attack pipelines. Meanwhile evaluation practice drifts the other way: newest defenses arrive with ever-cleaner galleries under ever-stronger laboratory attackers~\cite{fardin2026metacloakjpeg,chen2025fragile}, and the first cross-method studies still hold the acquisition side fixed~\cite{wang2025privacy}---a literature optimizing for an adversary that does not exist.

We build the missing instrument. \emph{ScrapeBench} reconstructs the funnel in software: seven tiers (T0 original; T1 platform transcode; T2 browser and thumbnail rendering; T3 agent-style capture; T4 super-resolution; T5 diffusion regeneration; T6 adaptive mixing and purification), each with an explicit attacker budget. On top of it we define the \emph{retention ratio} $\mathrm{RR}(T)=\Delta(T)/\Delta(T_0)$, where $\Delta(T)$ is the protection gap (clean minus protected identity similarity under a fixed attack) at tier $T$: RR is the fraction of a defense's T0 gap that survives acquisition reality, computed against a same-identity clean control. The primary evaluation applies each of six representative preprocessing operations (JPEG-85, WebP-80, browser render, agent face-crop, super-resolution, diffusion regeneration) \emph{in isolation}, because the marginal damage a single operation does to a protection gap is the quantity a defender can act on; stacking operations multiplies combinatorics faster than it multiplies insight, and a completed composition study (Appendix~\ref{sec:extra}) confirms that reordering or deepening the chain does not rescue a gap. The metric is paired with a cost ledger: the defender pays GPU-minutes per identity to forge a perturbation, the attacker pays CPU-milliseconds per image to defeat it.

Running ScrapeBench over 40 public identities, two published defenses (FSMG-style untargeted surrogate attacks~\cite{le2023antidreambooth} and ASPL-lite alternation), clean and random controls, a JPEG-aware retrained variant, an acquisition-robust reference defense, and three attack families (DreamBooth-LoRA fine-tuning and the zero-shot adapters IP-Adapter-FaceID and InstantID) yields five findings, each readable from the retention tables in Section~\ref{sec:results}: \textbf{(F1)} against same-identity clean controls the reported T0 gap shrinks roughly three-fold relative to the cross-identity baselines the field reports against, and loses significance under multiple-comparison correction; \textbf{(F2)} no published protection retains a statistically defensible gap under \emph{any} of the five signal-preserving single-tier operations---the sixth, diffusion regeneration, lifts every retention ratio above unity by degrading the photographs themselves---and a JPEG-aware retrain buys nothing at matched budget; \textbf{(F3)} what protection exists is \emph{reader-conditional}---large under fine-tuning, small under IP-Adapter-FaceID, sizable but underpowered under InstantID---so mechanism labels alone do not predict deployment behavior; \textbf{(F4)} the failure modes are specific enough to design against: ARCP, an acquisition-robust camouflage perturbation with a texture-allocated budget and a reader-diverse objective, is released as the reference defense the reporting standard calls for; \textbf{(F5)} the cost asymmetry runs against the defender by orders of magnitude---GPU-minutes of forging buy a gap that platform milliseconds narrow to noise. Two distinctions keep the interpretation honest: we separate \emph{defense bypass} from \emph{image destruction}, and we claim scope explicitly---two published $T_0$-optimized defenses, a JPEG-aware retrain, two controls, and an acquisition-robust reference defense, bounding deployment value rather than proving impossibility. The constructive turn is a funnel-aware, identity-matched reporting standard (RR@T against same-identity baselines), a reference defense designed for the acquisition distribution, and a consent layer that survives the funnel by construction.

Our contributions are:
\begin{itemize}[leftmargin=*,itemsep=1pt,topsep=2pt]
\item \textbf{The Web Acquisition Gap and a funnel threat model, anchored in platform reality}: a seven-tier funnel with per-tier attacker budgets, instrumented against public image services (GitHub avatars, Gravatar, WordPress Photon, an open proxy) whose measured re-encoding and resizing covers exactly the tiers that collapse the defenses.
\item \textbf{The retention-ratio metric and collapse-curve protocol}: $\mathrm{RR}(T)$ reduces robustness under acquisition to one interpretable curve per defense; a cost ledger exposes the defender--attacker asymmetry; an infeasibility view separates defense bypass from image destruction.
\item \textbf{ScrapeBench, an open harness}: transcoding, browser and agent rendering, restoration, regeneration, and adaptive pooling, integrated with DreamBooth-LoRA and two zero-shot attackers, ArcFace scoring, paired statistics, and automatic result publication.
\item \textbf{Empirical collapse findings with repaired pilots}: across 40 identities, two published defenses, two controls, a JPEG-aware variant, an acquisition-robust reference defense, and three attack families, quantifying F1--F5 over six single-tier preprocessing operations---establishing where, not just whether, protections fail.
\item \textbf{ARCP, an acquisition-robust reference defense}: a band-limited, texture-allocated perturbation optimized in expectation over differentiable acquisition proxies with a reader-diverse objective---the first protection whose training objective is the acquisition distribution rather than the T0 file, released under the benchmark's own reporting contract.
\item \textbf{A consent layer that survives the funnel}: \texttt{face-consent.txt}, a machine-readable consent manifest with a reference validator, shifting the durable line of defense from perturbations the funnel erases to acquisition-layer provenance the funnel carries with the image.
\end{itemize}
